Training runs in three steps on one label set, each step starting from the checkpoint of the step before; the patch grows from step to step and the batch shrinks after the first step. The checkpoint at the end of the first step is released as Rectina m1, and the checkpoint at the end of the third step as Rectina m2.

Stage 1 (s1), released as Rectina m1: 131000 iterations (taken at its fixed step, ckpt\_130001), batch 64 patches per step with per scroll quotas (0139 12, 1667 8, Paris4 6, 0814 7, 0009B 8, 0343P 5, 0500P2 7, 0172 11), patches of 17x128x128 voxels (z, y, x), learning rate 0.01, optimizer sgd, warmup 1000 steps, validation every 1000 steps, depth warp augmentation with a maximum offset of 3 planes. The base settings are those of the aligned21 hybrid 3D/2D configuration of the villa ink detection code (optimizer sgd, learning rate 0.01, batch size 64, patch 17x128x128, precision fp16). The run was continued in place, with the full training state and the data sampler advanced, from the checkpoint at iteration 79001 of an earlier run of the same stage.\footnote{The first training step was continued in place from iteration 79001 of an earlier run of the same recipe, itself trained from scratch.}

Stage 3 (s3), the first of the two further stages: from the s1 checkpoint ckpt\_130001, 30000 iterations, batch size 10, patches of 17x256x256 voxels, learning rate 0.003, validation every 2000 steps, the checkpoint of the last step (ckpt\_030000) is taken.

Stage 4 (s4), released as Rectina m2: from the s3 checkpoint of the last step, 20000 iterations, batch size 10, patches of 17x384x384 voxels, learning rate 0.003. The checkpoint released is the one of the last step (ckpt\_020000).

Two additions apply to stage 1 only. The first is the per patch depth warp, which replaces villa's flat random shift of the depth window by a smooth random shift of the planes read for each training patch: a smooth field built from Gaussian noise at scales of 8, 16, 32, 64 pixels, rounded to the nearest plane, with a maximum offset of 3 planes; empty rectangles of 32 pixels are blanked, with probability 0.05, in the image and in its supervision mask; labels are read without shift. The registration of the chain sets the range to 2, the maximum offset to 3 and the amplitude to 2. Stages 3 and 4 carry no depth warp.

The second is the robust\_mad guard of the training hook. Villa's robust\_mad normalisation divides a patch by its scaled median absolute deviation and falls back to a different scale only when that deviation is under 1e-6. A warped window that is about half fill (value 0) has a median at or near 0 and a tiny deviation, so the patch would be divided by a very small scale and reach huge values. The guard changes this for the image the warp has just built, and only in the robust\_mad mode: when the scaled deviation of the patch, after clipping to the 1st and 99th percentiles, is at least that floor but under 0.01 of the span between those percentiles, the patch is normalised with villa's own fallback branch instead of by the tiny scale; in every other case villa normalises as before. The guard can act only where a warp has just built the image, that is in stage 1.

Nothing else of the recipe is dropped or added after stage 1: the depth warp is dropped from stages 3 and 4, and the checkpoint of the last step of stage 4 is the one released as Rectina m2. Settings common to the three stages: mixed precision fp16, scheduler diffusers\_cosine\_warmup, weight decay 3e-05, label smoothing of the BCE loss 0.5, seed 42, 32 validation batches per evaluation.
